\documentclass[12pt,a4paper]{article}

\usepackage[utf8]{inputenc}
\usepackage[T1]{fontenc}
\usepackage[english]{babel}
\usepackage{amsmath}
\usepackage{geometry}
\usepackage{xcolor}
\usepackage{mdframed}
\usepackage{enumitem}

\geometry{margin=2.5cm}

\definecolor{azul}{RGB}{30, 100, 180}
\definecolor{cinza}{RGB}{245, 245, 245}
\definecolor{verde}{RGB}{20, 130, 80}

\mdfdefinestyle{enunciado}{
  backgroundcolor=cinza,
  linecolor=azul,
  linewidth=1.2pt,
  topline=false, rightline=false, bottomline=false,
  innerleftmargin=14pt, innerrightmargin=14pt,
  innertopmargin=12pt, innerbottommargin=12pt
}

\mdfdefinestyle{dados}{
  backgroundcolor=cinza,
  linecolor=gray!50,
  linewidth=0.5pt,
  innerleftmargin=12pt, innerrightmargin=12pt,
  innertopmargin=8pt, innerbottommargin=8pt,
  roundcorner=4pt
}

\mdfdefinestyle{resolucao}{
  backgroundcolor=white,
  linecolor=verde,
  linewidth=1.2pt,
  topline=false, rightline=false, bottomline=false,
  innerleftmargin=14pt, innerrightmargin=14pt,
  innertopmargin=12pt, innerbottommargin=12pt
}

\begin{document}

\begin{center}
  {\small\textcolor{azul}{\textbf{F\'isica e Qu\'imica A\,---\,11.\textordmasculine{} ano}}}
\end{center}

\bigskip

%---------------------------------------------------------------
\begin{mdframed}[style=enunciado]

Um investigador prepara uma solu\c{c}\~ao aquosa saturada de um sal pouco
sol\'uvel, XY, a 25\,\textdegree C, que se dissolve segundo o equil\'ibrio:
\[
  \mathrm{XY}_{(s)} \rightleftharpoons \mathrm{X}^{+}_{(aq)}
  + \mathrm{Y}^{-}_{(aq)}
\]
Para medir a concentra\c{c}\~ao da solu\c{c}\~ao, o investigador recorre a um
feixe de luz \textit{laser}: dirige-o do ar para a superf\'icie da solu\c{c}\~ao
com um \^angulo de incid\^encia de $45{,}0^\circ$ e mede um \^angulo de
refra\c{c}\~ao de $32{,}0^\circ$. Sabe-se que o \'indice de refra\c{c}\~ao de
uma solu\c{c}\~ao aquosa varia com a concentra\c{c}\~ao total de i\~oes segundo:
\[
  n_{\text{sol}} = n_{\text{\'{a}gua}} + k \cdot c_{\text{total}}
\]
Usando a lei de Snell-Descartes e a express\~ao acima, determine a
solubilidade molar de XY e verifique se o resultado \'e consistente com
o valor de $K_s$ fornecido.

\end{mdframed}

\medskip

%---------------------------------------------------------------
\begin{mdframed}[style=dados]
\textbf{Dados:}
\quad $K_s\,(\mathrm{XY},\;25\,\text{\textdegree C}) = 7{,}50 \times 10^{-4}$
\quad;\quad
$n_{\text{ar}} = 1{,}000$
\quad;\quad
$n_{\text{\'{a}gua}} = 1{,}333$
\quad;\quad
$k = 2{,}50 \times 10^{-2}\ \text{L\,mol}^{-1}$
\end{mdframed}

\bigskip\bigskip

%---------------------------------------------------------------
\noindent\textcolor{verde}{\rule{\linewidth}{0.5pt}}\\[4pt]
{\textcolor{verde}{\textbf{Resolu\c{c}\~ao}}}\\[-2pt]
\noindent\textcolor{verde}{\rule{\linewidth}{0.5pt}}

\bigskip

\begin{mdframed}[style=resolucao]

Pela lei de Snell-Descartes, o \'indice de refra\c{c}\~ao da solu\c{c}\~ao \'e:
\[
  n_{\text{sol}}
  = \frac{n_{\text{ar}}\,\sin 45{,}0^\circ}{\sin 32{,}0^\circ}
  = \frac{1{,}000 \times 0{,}7071}{0{,}5299}
  = 1{,}3344
\]

Substituindo na express\~ao do \'indice de refra\c{c}\~ao:
\[
  c_{\text{total}}
  = \frac{n_{\text{sol}} - n_{\text{\'{a}gua}}}{k}
  = \frac{1{,}3344 - 1{,}333}{2{,}50 \times 10^{-2}}
  = \frac{1{,}4 \times 10^{-3}}{2{,}50 \times 10^{-2}}
  = 5{,}47 \times 10^{-2}\ \text{mol\,L}^{-1}
\]

Como XY $\rightleftharpoons$ X$^+$ + Y$^-$, t\^em-se $[\mathrm{X}^+] =
[\mathrm{Y}^-] = s$ e $c_{\text{total}} = 2s$, logo:
\[
  s = \frac{c_{\text{total}}}{2}
    = \frac{5{,}47 \times 10^{-2}}{2}
    \approx 2{,}74 \times 10^{-2}\ \text{mol\,L}^{-1}
\]

\textit{Verifica\c{c}\~ao:} pelo $K_s$,
$s = \sqrt{K_s} = \sqrt{7{,}50 \times 10^{-4}} \approx 2{,}74 \times 10^{-2}$
mol\,L$^{-1}$. Os valores coincidem, confirmando a consist\^encia entre a
medida \'otica e o equil\'ibrio de solubilidade.

\end{mdframed}

\end{document}
